\documentclass[twocolumn,10pt,a4paper]{article}

\usepackage[utf8]{inputenc}
\usepackage[catalan]{babel}
\usepackage{graphicx}
\usepackage{amsmath}
\usepackage{geometry}
\usepackage{listings}
\usepackage{xcolor}
\usepackage{booktabs}
\usepackage{caption}

\geometry{margin=2.5cm, bottom=2.5cm}

\definecolor{codegreen}{rgb}{0,0.6,0}
\definecolor{codegray}{rgb}{0.5,0.5,0.5}
\definecolor{codepurple}{rgb}{0.58,0,0.82}
\definecolor{backcolour}{rgb}{0.95,0.95,0.92}

\lstset{
    backgroundcolor=\color{backcolour},   
    commentstyle=\color{codegreen},
    keywordstyle=\color{magenta},
    numberstyle=\tiny\color{codegray},
    stringstyle=\color{codepurple},
    basicstyle=\ttfamily\scriptsize,
    breakatwhitespace=false,         
    breaklines=true,                 
    captionpos=b,                    
    keepspaces=true,                 
    numbers=left,                    
    numbersep=5pt,                  
    showspaces=false,                
    showstringspaces=false,
    showtabs=false,                  
    tabsize=2
}

\title{\fontsize{16pt}{20pt}\selectfont \textbf{Arquitectura SCADA i Control de Tercer Nivell per a una Micro-xarxa DC Híbrida: Una Aproximació Ciber-Física}}

\author{
    Guillem Castillo i Arnau Coronado \\
    \small Escola d'Enginyeria de Barcelona Est (EEBE) \\
    \small Universitat Politècnica de Catalunya (UPC) \\
    \small \textit{Màster en Enginyeria Electrònica (PGE3)}
}
\date{1 de juny de 2026}

\begin{document}

\maketitle

\begin{abstract}
    Aquest informe documenta el disseny, implementació i validació d'un sistema de supervisió i control (SCADA) per a una micro-xarxa DC híbrida. El projecte integra una capa física basada en l'ESP32 amb un motor algorítmic en Python, permetent la validació de polítiques energètiques mitjançant tres modes d'execució progressius: simulació purament matemàtica, Hardware-in-the-Loop (HIL) i control de planta real. S'analitza amb profunditat el tractament de senyal en el microcontrolador mitjançant filtratge per mediana, la gestió de la comunicació asíncrona per minimitzar la latència i les estratègies de \textit{load shedding} i \textit{safety sink} implementades al mòdul de decisió.
\end{abstract}

\section{Introducció i Objectius}
Dins de la jerarquia de control de micro-xarxes, el tercer nivell s'ocupa de la gestió de potències i l'estabilitat a llarg termini del bus, optimitzant el balanç entre generació, demanda i emmagatzematge. L'objectiu d'aquest treball és desenvolupar una plataforma ciber-física funcional que permeti emular i controlar aquests fluxos d'energia.

Els reptes tecnològics abordats inclouen:
\begin{itemize}
    \item Disseny d'un motor de simulació basat en el mètode d'Euler.
    \item Implementació de Hardware-in-the-Loop per validar la latència del protocol sèrie.
    \item Robustesa en l'adquisició davant el soroll de commutació PWM a 1kHz.
    \item Estabilització del bus a 40V mitjançant control proporcional de la xarxa elèctrica.
\end{itemize}

\section{Arquitectura del Sistema}
La infraestructura es divideix en components de potència i etapes de control:
\begin{enumerate}
    \item \textbf{Generació Emulada:} S'utilitzen les sortides DAC de l'ESP32 per controlar fonts de corrent que emulen panells FV i la xarxa elèctrica.
    \item \textbf{Emmagatzematge:} Un bus DC amb un condensador de 4700$\mu$F que actua com a memòria d'energia.
    \item \textbf{Branques de Càrrega:} Tres LEDs controlats per transistors BJT/MOSFET en configuració d'emissor comú.
\end{enumerate}

\subsection{Mapeig de Pins (ESP32)}
\begin{itemize}
    \item \textbf{ADC 34 (Bus):} Entrada crítica que alimenta la lògica del \textit{Brain}.
    \item \textbf{DAC 25/26:} Control de generació. En mode HIL, el DAC 26 injecta el voltatge simulat.
    \item \textbf{PWM 2, 4, 5:} Control de \textit{duty cycle} (0-255) per a les càrregues Roig, Blau i Verd.
\end{itemize}

\section{Ecosistema de Programari}
La complexitat del projecte resideix en la seva modularitat. Cada script té una responsabilitat única dins del llaç de control de 50ms.

\subsection{Firmware i Filtrat Digital}
L'ESP32 no es limita a ser un transceptor de dades. Realitza un preprocessament digital essencial. El soroll induït pel PWM en un bus compartit fa que les lectures ADC siguin inestables. Per solucionar-ho, \texttt{uControl\_0.py} implementa un filtratge per mediana:

\begin{lstlisting}[language=Python]
def read_stable(adc):
    samples = []
    for _ in range(15):
        samples.append(adc.read())
    samples.sort()
    return samples[7] 
\end{lstlisting}
Aquesta tècnica descarta els valors atípics (*outliers*) produïts pels transitoris de commutació, oferint una mesura de tensió de bus extremadament neta al SCADA.

\subsection{Abstracció de Hardware}
El mòdul \texttt{control\_esp32.py} resol el problema del retard en Python. Els sistemes SCADA sovint pateixen de \textit{lag} si el buffer sèrie s'omple. La nostra implementació buida el buffer i busca l'última línia completa des del final:

\begin{lstlisting}[language=Python]
if self.ser.in_waiting > 0:
    raw_data = self.ser.read(self.ser.in_waiting).decode()
    lines = raw_data.strip().split('\n')
    latest_line = lines[-1] 
\end{lstlisting}

\subsection{Model de Física Digital}
El mòdul \texttt{simulation.py} permet al sistema predir el comportament del bus. Utilitzant les lleis de Kirchhoff i el mètode d'Euler, s'actualitza l'estat del condensador:
\begin{equation}
    V_{bus_{k+1}} = V_{bus_k} + \frac{I_{solar} + I_{grid} - I_{cons} - I_{loss}}{C} \cdot \Delta t
\end{equation}
S'ha modelat una pèrdua $I_{loss} = V_{bus}/200$ per evitar que el sistema sigui un integrador pur ideal, emulant així les fuites reals de qualsevol sistema d'emmagatzematge.

\section{Modes d'Operació}
L'habilitat del projecte per commutar entre modes sense canviar el \textit{Brain} és un indicador de bon disseny de programari:

\begin{itemize}
    \item \textbf{Simulació Pura:} Validació teòrica.
    \item \textbf{Mode Híbrid (HIL):} \texttt{hybrid.py} tanca el bucle DAC-ADC. És el mode més exigent per al firmware, ja que l'ESP32 ha de digitalitzar un senyal que ell mateix ha generat després d'un processament al PC.
    \item \textbf{Mode Hardware Real:} \texttt{main.py} s'executa sobre la planta física. Les dades ADC (Pin 34) són ara l'única veritat del sistema, eliminant la dependència del simulador.
\end{itemize}

\section{Estratègies de Control i Seguretat}
El mòdul \texttt{brain.py} actua com l'òrgan de decisió central, aplicant les següents consignes:

\subsection{Control Proporcional del Grid}
Per mantenir el bus a la consigna de 40V, s'injecta potència de xarxa de forma proporcional a l'error de tensió:
\begin{itemize}
    \item \textbf{$KP_{GRID} = 500$:} Un guany elevat per a una resposta ràpida.
    \item \textbf{Grid Throttle:} Una protecció que disminueix la injecció de xarxa si la tensió supera els 40V, evitant sobreoscil·lacions i estalviant energia.
\end{itemize}

\subsection{Jerarquia de Prioritats i Load Shedding}
La micro-xarxa prioritza la continuïtat del servei de la càrrega crítica (Vermella). El \textit{load shedding} actua de forma seqüencial:
\begin{enumerate}
    \item \textbf{Zona Operativa ($>38$V):} Totes les càrregues actives.
    \item \textbf{Llindar 1 (38V):} Es desconnecta la càrrega Verda.
    \item \textbf{Llindar 2 (36V):} Es desconnecta la càrrega Blava.
    \item \textbf{Reserva Crítica:} El LED Vermell només s'apaga si la tensió del bus cau de 30V (límit físic de seguretat).
\end{enumerate}

Protecció Avançada: Safety Sink
Si la generació solar és màxima i no hi ha consum, la tensió del bus pot pujar ràpidament. A 45V, el \textit{Brain} activa el \textbf{Safety Sink}: força el LED verd al 100\% independentment del slider de l'usuari per drenar energia. Si la tensió persisteix per sobre de 45V, el DAC solar es posa a 0 per tallar la producció.

\section{Anàlisi i Resultats}
La integració del fitxer \texttt{irradiancia.csv} ha permès realitzar assajos de llarga durada (cicles de 24h emulats). Els gràfics de \textit{pyqtgraph} mostren una estabilitat del bus amb un error inferior a l'1\% en mode hardware real. La resposta del \textit{load shedding} és gairebé instantània (una iteració del bucle, 50ms), el que garanteix la supervivència de la càrrega roja fins i tot en caigudes sobtades de producció solar.

\section{Conclusions}
Aquest projecte demostra que un control de tercer nivell robust depèn tant de la qualitat de l'algorisme com de la integritat de l'adquisició de dades. El filtrat digital a l'ESP32 i la gestió de buffers a Python han estat tan decisius per a l'estabilitat com la pròpia lògica de control proporcional. L'arquitectura modular aconseguida és directament escalable a micro-xarxes de major potència.

\end{document}
